\documentclass[aspectratio=169]{beamer}

% Theme and color scheme
\usetheme{default}
\usecolortheme{default}

% Remove navigation symbols
\setbeamertemplate{navigation symbols}{}

% Simple color scheme for printing
\definecolor{darkblue}{RGB}{25,25,112}
\setbeamercolor{title}{fg=darkblue}
\setbeamercolor{frametitle}{fg=darkblue}
\setbeamercolor{structure}{fg=darkblue}

% Simple footline
\setbeamertemplate{footline}[frame number]

% Packages
\usepackage{amsmath}
\usepackage{amsfonts}
\usepackage{amssymb}
\usepackage{graphicx}
\usepackage{tikz}
\usepackage{booktabs}
\usepackage{array}
\usepackage{listings}
\usepackage{xcolor}

% Set graphics path for images (relative to slides directory)
\graphicspath{{../Images/}}

% Title information
\title[Chapter 3.14: Parallel Computing]{Structured Computer Architecture\\Chapter 3.14: Parallel Computing\\N-Order Systems (N-OS)}
\author{Based on Structured Computer Architecture: An Order-Based Approach}
\date{}

% Custom command for highlighting concepts
\newcommand{\concept}[1]{\textbf{\color{darkblue}#1}}

\begin{document}

% Title slide
\begin{frame}
\titlepage
\end{frame}

% Reset section counter for this chapter
\setcounter{section}{0}

% Define sections for this chapter
\section{Introduction to N-Order Systems}
\section{Cellular Automata Theory}
\section{Controlled Cellular Automata}
\section{Parallel Computing Architectures}
\section{Mapping and Distribution}
\section{Performance and Scalability}

% Table of contents
\begin{frame}{Outline}
\tableofcontents
\end{frame}

% Learning objectives
\begin{frame}{Learning Objectives}
By the end of this chapter, you will understand:
\begin{itemize}
\item N-Order Systems (N-OS) and their theoretical foundations
\item Cellular Automata concepts and applications
\item Controlled Cellular Automata (CCA) architecture
\item Parallel computing paradigms and mapping operations
\item Distribution networks and scalability challenges
\item Space-time trade-offs in parallel computation
\item Moore's Law implications for parallel systems
\end{itemize}
\end{frame}

% Section 1: Introduction to N-Order Systems
\section{Introduction to N-Order Systems}

\begin{frame}{From Sequential to Parallel: The N-OS Evolution}
\begin{columns}
\begin{column}{0.6\textwidth}
\concept{Motivation for N-Order Systems}:
\begin{itemize}
\item Previous chapters: Sequential feedback loops (1-OS to 5-OS)
\item Question: Can we continue closing new loops?
\item Goal: Additional gains from additional loops
\item Challenge: Managing complexity as N increases
\end{itemize}

\vspace{0.3cm}
\concept{Complexity Management}:
\begin{itemize}
\item \textbf{Size}: Proportional to N (acceptable)
\item \textbf{Complexity}: Must remain independent of N
\item \textbf{Solution}: Loops close over identical systems
\item \textbf{Repetitive patterns}: Keep description controllable
\end{itemize}

\vspace{0.3cm}
\concept{Moore's Law Constraint}:
\begin{itemize}
\item Exponential growth in size (transistors)
\item Complexity must remain manageable
\item Repetitive patterns enable scaling
\end{itemize}
\end{column}
\begin{column}{0.4\textwidth}
\begin{center}
\textbf{System Order Evolution}
\vspace{0.3cm}

\tikz{
% System progression
\node at (1,4) {\textbf{System Evolution}};

\draw[thick] (0,3.5) rectangle (2,3.7);
\node at (1,3.6) {0-OS: Combinational};

\draw[thick] (0,3.1) rectangle (2,3.3);
\node at (1,3.2) {1-OS: Memory};

\draw[thick] (0,2.7) rectangle (2,2.9);
\node at (1,2.8) {2-OS: Automata};

\draw[thick] (0,2.3) rectangle (2,2.5);
\node at (1,2.4) {3-OS: Processors};

\draw[thick] (0,1.9) rectangle (2,2.1);
\node at (1,2) {4-OS: Von Neumann};

\draw[thick] (0,1.5) rectangle (2,1.7);
\node at (1,1.6) {5-OS: Harvard};

\draw[thick, dashed] (0,1.1) rectangle (2,1.3);
\node at (1,1.2) {N-OS: Parallel};

% Arrows
\foreach \y in {3.3,2.9,2.5,2.1,1.7} {
  \draw[->] (1,\y) -- (1,\y-0.2);
}
}
\end{center}
\end{column}
\end{columns}
\end{frame}

\begin{frame}{N-Order System Structure}
\begin{columns}
\begin{column}{0.6\textwidth}
\concept{N-OS Characteristics}:
\begin{itemize}
\item \textbf{Size}: $O(N)$ - Linear growth with number of components
\item \textbf{Complexity}: $O(1)$ or $O(N)$ depending on design
\item \textbf{Simple Systems}: Identical cells, $O(1)$ complexity
\item \textbf{Complex Systems}: Different cells, $O(N)$ complexity
\end{itemize}

\vspace{0.3cm}
\concept{Design Requirements}:
\begin{itemize}
\item Cells $c_0, c_1, \ldots, c_{N-1}$ must be identical
\item Regular interconnection pattern
\item Description involves:
\begin{itemize}
\item Single cell description
\item Interconnection pattern for N-1 cells
\end{itemize}
\end{itemize}

\vspace{0.3cm}
\concept{Scalability Principle}:
\begin{itemize}
\item Only simple systems (identical cells) are practical
\item Complex systems become uncontrollable for large N
\item Goal: Maximize N while maintaining simplicity
\end{itemize>
\end{column>
\begin{column>{0.4\textwidth>
\begin{center>
\textbf{N-OS Structure}
\vspace{0.3cm}

\tikz{
% N-OS representation
\node at (1,3.5) {\textbf{N-Order System}};

% Cells
\foreach \x/\label in {0/c_0, 0.8/c_1, 1.6/c_2, 2.4/c_{N-1}} {
  \draw[thick] (\x,2.5) rectangle (\x+0.6,3);
  \node at (\x+0.3,2.75) {\label};
}

% Connections
\draw[<->] (0.6,2.75) -- (0.8,2.75);
\draw[<->] (1.4,2.75) -- (1.6,2.75);
\draw[dashed] (2.2,2.75) -- (2.4,2.75);

% Labels
\node at (1.5,2.2) {Identical Cells};
\node at (1.5,1.9) {Regular Connections};

% Complexity comparison
\node at (1.5,1.4) {\textbf{Complexity}};
\node at (0.7,1.1) {Simple: $O(1)$};
\node at (2.3,1.1) {Complex: $O(N)$};
\draw[thick, green] (0.2,1.1) -- (1.2,1.1);
\draw[thick, red] (1.8,1.1) -- (2.8,1.1);
}
\end{center>
\end{column>
\end{columns>
\end{frame>

% Section 2: Cellular Automata Theory
\section{Cellular Automata Theory}

\begin{frame}{Theoretical Cellular Automaton}
\begin{columns}
\begin{column}{0.6\textwidth}
\concept{Cellular Automaton Definition}:
\begin{itemize}
\item Unlimited number of identical cells
\item Linear connection between adjacent cells
\item Each cell is a finite automaton
\item Synchronous operation across all cells
\end{itemize}

\vspace{0.3cm}
\concept{Cell Behavior}:
\begin{itemize}
\item Internal state determines cell behavior
\item Inputs from left and right neighbors
\item State transitions based on:
\begin{itemize}
\item Current internal state
\item Left neighbor signal
\item Right neighbor signal
\end{itemize>
\item All cells switch simultaneously
\end{itemize>

\vspace{0.3cm}
\concept{Theoretical Properties}:
\begin{itemize}
\item Infinite array of cells
\item Homogeneous structure
\item Local interactions only
\item Global behavior emerges from local rules
\end{itemize>
\end{column>
\begin{column>{0.4\textwidth>
\begin{center>
\textbf{Cellular Automaton}
\vspace{0.3cm}

\tikz{
% Cellular automaton representation
\node at (1.5,3.5) {\textbf{1D Cellular Automaton}};

% Cells
\foreach \x/\label in {0/..., 0.8/c_{i-1}, 1.6/c_i, 2.4/c_{i+1}, 3.2/...} {
  \draw[thick] (\x,2.5) rectangle (\x+0.6,3);
  \node at (\x+0.3,2.75) {\label};
}

% Connections
\draw[<->] (0.6,2.75) -- (0.8,2.75);
\draw[<->] (1.4,2.75) -- (1.6,2.75);
\draw[<->] (2.2,2.75) -- (2.4,2.75);
\draw[<->] (3,2.75) -- (3.2,2.75);

% State transitions
\draw[->] (1.9,2.5) -- (1.9,2);
\node at (1.9,1.8) {State};
\node at (1.9,1.6) {Transition};

% Input signals
\draw[->] (1.4,2.3) -- (1.6,2.3);
\draw[->] (2.2,2.3) -- (2,2.3);
\node at (1.2,2.1) {Left};
\node at (2.4,2.1) {Right};

% Clock
\node at (1.5,1.2) {Synchronous Clock};
}
\end{center>
\end{column>
\end{columns>
\end{frame>

\begin{frame}{Simple Cellular Automaton Example}
\begin{columns}
\begin{column}{0.6\textwidth}
\concept{One-Bit Cell Automaton}:
\begin{itemize}
\item Each cell has one-bit state
\item Two one-bit inputs (left and right neighbors)
\item State transition function based on 3-bit input
\item 8 possible input combinations
\end{itemize>

\vspace{0.3cm}
\concept{Implementation Characteristics}:
\begin{itemize}
\item Two-state finite automaton per cell
\item General form using 8-bit lookup table
\item Multiplexer-based implementation
\item Constant description size (independent of N)
\end{itemize>

\vspace{0.3cm}
\concept{Key Properties}:
\begin{itemize}
\item \textbf{Simplicity}: Constant complexity $O(1)$
\item \textbf{Scalability}: Size independent description
\item \textbf{Regularity}: Identical cells and connections
\item \textbf{Efficiency}: Parallel operation
\end{itemize>
\end{column>
\begin{column>{0.4\textwidth>
\begin{center>
\textbf{Cell State Transition}
\vspace{0.3cm}

\tikz{
% Cell representation
\draw[thick] (0.5,2.5) rectangle (2.5,3.5);
\node at (1.5,3) {Cell $c_i$};

% Inputs
\draw[->] (0,3) -- (0.5,3);
\node at (0.2,3.3) {Left};
\draw[->] (3,3) -- (2.5,3);
\node at (2.8,3.3) {Right};

% State
\draw[thick] (1.2,2.5) rectangle (1.8,2.8);
\node at (1.5,2.65) {State};

% Transition function
\node at (1.5,2.2) {Transition Function};
\node at (1.5,1.9) {$f(state, left, right)$};

% Truth table representation
\node at (1.5,1.5) {\textbf{8-bit Lookup}};
\draw[thick] (0.8,1.1) rectangle (2.2,1.3);
\node at (1.5,1.2) {Rule Table};

% Output
\draw[->] (1.5,1.1) -- (1.5,0.8);
\node at (1.5,0.6) {Next State};
}
\end{center>
\end{column>
\end{columns>
\end{frame>

% Section 3: Controlled Cellular Automata
\section{Controlled Cellular Automata}

\begin{frame}{Limitations of Theoretical Model}
\begin{columns}
\begin{column}{0.6\textwidth}
\concept{Theoretical Model Problems}:
\begin{itemize}
\item \textbf{Computational Inefficiency}: Real problems cannot be solved effectively
\item \textbf{Unstructured Mechanism}: Completely unstructured computation
\item \textbf{Hardware Size}: Physical structure becomes inefficiently large
\item \textbf{All-or-Nothing}: All components involved in every cycle
\end{itemize>

\vspace{0.3cm}
\concept{Need for Practical Model}:
\begin{itemize}
\item Space-time trade-off capability
\item Suitable sequencing of computation steps
\item Structured state space
\item Controlled execution flow
\end{itemize>

\vspace{0.3cm}
\concept{Solution Approach}:
\begin{itemize}
\item Move from theoretical to abstract model
\item Increase flexibility of computing process
\item Reduce hardware structure size
\item Maintain parallel processing benefits
\end{itemize>
\end{column>
\begin{column>{0.4\textwidth>
\begin{center>
\textbf{Model Evolution}
\vspace{0.3cm}

\tikz{
% Theoretical model
\node at (1,3.5) {\textbf{Theoretical CA}};
\draw[thick] (0,3) rectangle (2,3.3);
\node at (1,3.15) {Unstructured};
\draw[thick] (0,2.7) rectangle (2,3);
\node at (1,2.85) {All Components};
\draw[thick] (0,2.4) rectangle (2,2.7);
\node at (1,2.55) {Every Cycle};

% Arrow
\draw[->, thick] (1,2.2) -- (1,1.8);
\node at (1.5,2) {Improve};

% Practical model
\node at (1,1.5) {\textbf{Controlled CA}};
\draw[thick] (0,1) rectangle (2,1.3);
\node at (1,1.15) {Structured};
\draw[thick] (0,0.7) rectangle (2,1);
\node at (1,0.85) {Selective Use};
\draw[thick] (0,0.4) rectangle (2,0.7);
\node at (1,0.55) {Sequenced};
}
\end{center>
\end{column>
\end{columns>
\end{frame>

\begin{frame}{Controlled Cellular Automaton (CCA) Architecture}
\begin{columns}
\begin{column}{0.6\textwidth}
\concept{CCA Components}:
\begin{itemize}
\item \textbf{SQ (Controller)}: Processor + data memory + program memory
\item \textbf{Distribute Network}: Log-depth distribution network
\item \textbf{Map (Cell Array)}: Finite number of cells with RALU units
\end{itemize>

\vspace{0.3cm}
\concept{Enhanced Cell Structure}:
\begin{itemize>
\item Replace simple state register with register file
\item Structured state space: $Q = \{Q_0 \times Q_1 \times \ldots \times Q_{m-1}\}$
\item RALU-type execution unit in each cell
\item Local memory resources
\end{itemize>

\vspace{0.3cm}
\concept{Control Mechanism}:
\begin{itemize}
\item SQ issues same command to all cells each cycle
\item Commands include: operations, addressing modes, addresses, constants
\item Serial connection for data transfer between SQ and cells
\item Distributed execution with centralized control
\end{itemize>
\end{column>
\begin{column>{0.4\textwidth>
\begin{center>
\textbf{CCA Architecture}
\vspace{0.3cm}

\tikz{
% Controller
\draw[thick] (0.5,3.5) rectangle (2.5,4);
\node at (1.5,3.75) {SQ Controller};

% Distribution network
\draw[thick] (1.2,3) rectangle (1.8,3.3);
\node at (1.5,3.15) {Dist};
\draw[->] (1.5,3.5) -- (1.5,3.3);

% Distribution tree
\draw[->] (1.5,3) -- (1,2.7);
\draw[->] (1.5,3) -- (2,2.7);
\draw[->] (1,2.7) -- (0.7,2.4);
\draw[->] (1,2.7) -- (1.3,2.4);
\draw[->] (2,2.7) -- (1.7,2.4);
\draw[->] (2,2.7) -- (2.3,2.4);

% Cells
\foreach \x/\label in {0.5/c_0, 1.1/c_1, 1.7/c_2, 2.3/c_3} {
  \draw[thick] (\x,2) rectangle (\x+0.4,2.4);
  \node at (\x+0.2,2.2) {\label};
}

% Cell connections
\draw[<->] (0.9,2.2) -- (1.1,2.2);
\draw[<->] (1.5,2.2) -- (1.7,2.2);

% Labels
\node at (1.5,1.7) {Map Array};
\node at (1.5,1.4) {RALU Cells};
}
\end{center>
\end{column>
\end{columns>
\end{frame>

\begin{frame}{Distribution Network Design}
\begin{columns}
\begin{column}{0.6\textwidth}
\concept{Distribution Network Requirements}:
\begin{itemize}
\item \textbf{Problem}: Direct connection from SQ to many cells impossible
\item \textbf{Electrical Constraints}: High frequency, large distances
\item \textbf{Solution}: Log-depth distribution tree
\item \textbf{Latency}: Logarithmic in number of cells
\end{itemize}

\vspace{0.3cm}
\concept{Network Characteristics}:
\begin{itemize}
\item Pipeline registers at each node
\item Two connections per node maximum
\item Reasonable distances between connections
\item Scalable to large number of cells
\end{itemize>

\vspace{0.3cm}
\concept{Performance Benefits}:
\begin{itemize}
\item Maintains high clock frequency
\item Enables large-scale parallel systems
\item Predictable latency characteristics
\item Efficient command distribution
\end{itemize>
\end{column>
\begin{column>{0.4\textwidth>
\begin{center>
\textbf{Distribution Network}
\vspace{0.3cm}

\tikz{
% Root
\draw[thick] (1.5,3.5) rectangle (1.8,3.8);
\node at (1.65,3.65) {SQ};

% Level 1
\draw[thick] (0.8,3) rectangle (1.1,3.3);
\draw[thick] (2.2,3) rectangle (2.5,3.3);
\draw[->] (1.65,3.5) -- (0.95,3.3);
\draw[->] (1.65,3.5) -- (2.35,3.3);

% Level 2
\draw[thick] (0.3,2.5) rectangle (0.6,2.8);
\draw[thick] (1.3,2.5) rectangle (1.6,2.8);
\draw[thick] (1.7,2.5) rectangle (2,2.8);
\draw[thick] (2.7,2.5) rectangle (3,2.8);

\draw[->] (0.95,3) -- (0.45,2.8);
\draw[->] (0.95,3) -- (1.45,2.8);
\draw[->] (2.35,3) -- (1.85,2.8);
\draw[->] (2.35,3) -- (2.85,2.8);

% Cells
\foreach \x in {0.45, 1.45, 1.85, 2.85} {
  \draw[thick] (\x-0.1,2) rectangle (\x+0.1,2.3);
  \draw[->] (\x,2.5) -- (\x,2.3);
}

% Labels
\node at (1.65,1.7) {Log-depth: $O(\log p)$};
\node at (1.65,1.4) {Pipeline registers};
\node at (1.65,1.1) {High frequency};
}
\end{center>
\end{column>
\end{columns>
\end{frame>

% Section 4: Parallel Computing Architectures
\section{Parallel Computing Architectures}

\begin{frame}{Mapping Operations in Parallel Computing}
\begin{columns}
\begin{column}{0.6\textwidth}
\concept{Mapping Concept}:
\begin{itemize}
\item Apply same function to each element in collection
\item Fundamental parallel computing operation
\item Examples: vector operations, array processing
\item Natural fit for cellular automaton structure
\end{itemize}

\vspace{0.3cm}
\concept{Implementation in CCA}:
\begin{itemize>
\item Map array contains cells for parallel execution
\item Same function applied to distributed data
\item Vector elements distributed across cells
\item Simultaneous execution on all elements
\end{itemize>

\vspace{0.3cm}
\concept{Programming Model}:
\begin{itemize}
\item High-level: \texttt{map(function, vector)}
\item Hardware: Distributed function execution
\item Control: Single instruction, multiple data (SIMD)
\item Efficiency: Parallel processing of large datasets
\end{itemize>
\end{column>
\begin{column>{0.4\textwidth>
\begin{center>
\textbf{Mapping Example}
\vspace{0.3cm}

\tikz{
% Input vector
\node at (1.5,3.5) {\textbf{Vector Doubling}};
\node at (0.5,3.2) {Input:};
\foreach \x/\val in {1/1, 1.4/2, 1.8/3, 2.2/4} {
  \draw[thick] (\x,3) rectangle (\x+0.3,3.3);
  \node at (\x+0.15,3.15) {\val};
}

% Function
\node at (1.5,2.7) {Function: $f(x) = 2x$};

% Parallel application
\foreach \x in {1.15, 1.55, 1.95, 2.35} {
  \draw[->] (\x,3) -- (\x,2.4);
  \node at (\x,2.5) {$\times 2$};
}

% Output vector
\node at (0.5,2.1) {Output:};
\foreach \x/\val in {1/2, 1.4/4, 1.8/6, 2.2/8} {
  \draw[thick] (\x,1.9) rectangle (\x+0.3,2.2);
  \node at (\x+0.15,2.05) {\val};
}

% Parallel execution
\node at (1.5,1.6) {Parallel Execution};
\node at (1.5,1.3) {All elements simultaneously};
}
\end{center>
\end{column>
\end{columns>
\end{frame>

\begin{frame}{CCA Internal State Representation}
\begin{columns}
\begin{column}{0.6\textwidth}
\concept{Matrix State Model}:
\begin{itemize}
\item Global state represented as matrix $\mathbf{M_{mp}}$
\item $m$ rows: memory locations per cell
\item $p$ columns: number of cells
\item Each element: data word in cell's local memory
\end{itemize>

\vspace{0.3cm}
\concept{Memory Organization}:
\begin{itemize}
\item Each cell has local memory with $m$ locations
\item Distributed memory across $p$ cells
\item Total memory: $m \times p$ words
\item Parallel access to different memory locations
\end{itemize>

\vspace{0.3cm}
\concept{Data Distribution}:
\begin{itemize}
\item Vectors distributed row-wise or column-wise
\item Matrices distributed across multiple cells
\item Load balancing through proper distribution
\item Communication between adjacent cells
\end{itemize>
\end{column>
\begin{column>{0.4\textwidth>
\begin{center>
\textbf{Matrix State Model}
\vspace{0.3cm}

\tikz{
% Matrix representation
\node at (1.5,3.5) {$\mathbf{M_{mp}}$};

% Matrix
\draw[thick] (0.5,2.5) rectangle (2.5,3.2);

% Grid lines
\foreach \x in {0.9, 1.3, 1.7, 2.1} {
  \draw (\x,2.5) -- (\x,3.2);
}
\foreach \y in {2.7, 2.9, 3.1} {
  \draw (0.5,\y) -- (2.5,\y);
}

% Labels
\node at (0.3,2.85) {$m$};
\node at (1.5,2.3) {$p$ cells};

% Cell indices
\foreach \x/\idx in {0.7/0, 1.1/1, 1.5/2, 1.9/..., 2.3/p-1} {
  \node at (\x,2.1) {$c_{\idx}$};
}

% Memory locations
\foreach \y/\idx in {3.05/0, 2.85/1, 2.65/..., 2.45/m-1} {
  \node at (0.2,\y) {$\idx$};
}

% Data distribution
\node at (1.5,1.8) {\textbf{Data Distribution}};
\node at (1.5,1.5) {Row: Vector elements};
\node at (1.5,1.2) {Column: Cell memory};
}
\end{center>
\end{column>
\end{columns>
\end{frame>

% Section 5: Mapping and Distribution
\section{Mapping and Distribution}

\begin{frame}{Instruction Format and Execution}
\begin{columns}
\begin{column}{0.6\textwidth}
\concept{Instruction Components}:
\begin{itemize}
\item \textbf{Logical-Arithmetic Operations}: Function to be performed
\item \textbf{Addressing Modes}: How to access local memory
\item \textbf{Addresses}: Memory locations in each cell
\item \textbf{Constants}: Immediate values for operations
\end{itemize}

\vspace{0.3cm}
\concept{Execution Model}:
\begin{itemize}
\item Single instruction broadcast to all cells
\item Each cell executes same operation
\item Different data in each cell's local memory
\item SIMD (Single Instruction, Multiple Data) paradigm
\end{itemize>

\vspace{0.3cm}
\concept{Data Transfer}:
\begin{itemize}
\item Serial connection between SQ and cells
\item Data movement between SQ memory and cell memory
\item Efficient bulk data transfer mechanisms
\item Support for various data distribution patterns
\end{itemize>
\end{column>
\begin{column>{0.4\textwidth>
\begin{center>
\textbf{SIMD Execution}
\vspace{0.3cm}

\tikz{
% Controller instruction
\draw[thick] (0.5,3.5) rectangle (2.5,3.8);
\node at (1.5,3.65) {SQ: ADD R1, R2, R3};

% Distribution
\draw[->] (1.5,3.5) -- (0.7,3);
\draw[->] (1.5,3.5) -- (1.5,3);
\draw[->] (1.5,3.5) -- (2.3,3);

% Cells executing same instruction
\foreach \x/\label in {0.5/c_0, 1.3/c_1, 2.1/c_2} {
  \draw[thick] (\x,2.7) rectangle (\x+0.4,3);
  \node at (\x+0.2,2.85) {\label};
}

% Different data
\node at (0.7,2.5) {Data A};
\node at (1.5,2.5) {Data B};
\node at (2.3,2.5) {Data C};

% Same operation
\node at (0.7,2.2) {A+A};
\node at (1.5,2.2) {B+B};
\node at (2.3,2.2) {C+C};

% Results
\node at (0.7,1.9) {2A};
\node at (1.5,1.9) {2B};
\node at (2.3,1.9) {2C};

\node at (1.5,1.5) {\textbf{Same Instruction}};
\node at (1.5,1.2) {\textbf{Different Data}};
}
\end{center>
\end{column>
\end{columns>
\end{frame>

\begin{frame}{Space-Time Trade-offs}
\begin{columns}
\begin{column}{0.6\textwidth}
\concept{Trade-off Principles}:
\begin{itemize}
\item \textbf{More Space}: More parallel units, faster execution
\item \textbf{Less Space}: Fewer units, sequential execution
\item \textbf{Optimal Balance}: Match problem size to hardware
\item \textbf{Reconfigurable}: Adapt to different problem sizes
\end{itemize}

\vspace{0.3cm}
\concept{Implementation Strategies}:
\begin{itemize}
\item \textbf{Temporal Sequencing}: Reuse hardware over time
\item \textbf{Spatial Distribution}: Distribute computation across space
\item \textbf{Hybrid Approaches}: Combine temporal and spatial methods
\item \textbf{Dynamic Allocation}: Adjust resources based on workload
\end{itemize>

\vspace{0.3cm}
\concept{Performance Considerations}:
\begin{itemize}
\item Communication overhead vs. computation
\item Memory bandwidth limitations
\item Synchronization requirements
\item Load balancing across cells
\end{itemize>
\end{column>
\begin{column>{0.4\textwidth>
\begin{center>
\textbf{Space-Time Trade-off}
\vspace{0.3cm}

\tikz{
% Axes
\draw[->] (0,0) -- (3,0);
\draw[->] (0,0) -- (0,3);
\node at (3.2,0) {Space};
\node at (0,3.2) {Time};

% Trade-off curve
\draw[thick, blue] (0.5,2.5) .. controls (1.5,1.5) .. (2.5,0.5);

% Points
\fill[red] (0.7,2.2) circle (0.05);
\node at (0.7,2.5) {Sequential};
\fill[red] (2.3,0.7) circle (0.05);
\node at (2.3,0.4) {Parallel};

% Optimal region
\fill[green] (1.5,1.5) circle (0.05);
\node at (1.5,1.8) {Optimal};

% Labels
\node at (1.5,3.5) {\textbf{Trade-off Curve}};
}
\end{center>
\end{column>
\end{columns>
\end{frame>

% Section 6: Performance and Scalability
\section{Performance and Scalability}

\begin{frame}{Scalability Analysis}
\begin{columns}
\begin{column}{0.6\textwidth}
\concept{Scalability Metrics}:
\begin{itemize}
\item \textbf{Strong Scaling}: Fixed problem size, increase processors
\item \textbf{Weak Scaling}: Problem size grows with processors
\item \textbf{Efficiency}: Actual speedup vs. theoretical maximum
\item \textbf{Overhead}: Communication and synchronization costs
\end{itemize>

\vspace{0.3cm}
\concept{Limiting Factors}:
\begin{itemize}
\item \textbf{Amdahl's Law}: Sequential portions limit speedup
\item \textbf{Communication Latency}: Distribution network delays
\item \textbf{Memory Bandwidth}: Data transfer bottlenecks
\item \textbf{Load Imbalance}: Uneven work distribution
\end{itemize>

\vspace{0.3cm}
\concept{Optimization Strategies}:
\begin{itemize}
\item Minimize communication overhead
\item Optimize data locality
\item Balance computational load
\item Overlap computation and communication
\end{itemize>
\end{column>
\begin{column>{0.4\textwidth>
\begin{center>
\textbf{Scalability Curves}
\vspace{0.3cm}

\tikz{
% Axes
\draw[->] (0,0) -- (3,0);
\draw[->] (0,0) -- (0,3);
\node at (3.2,0) {Processors};
\node at (0,3.2) {Speedup};

% Ideal scaling
\draw[thick, green] (0,0) -- (2.8,2.8);
\node at (2.5,2.5) {Ideal};

% Realistic scaling
\draw[thick, blue] (0,0) .. controls (1,0.8) and (2,1.5) .. (2.8,1.8);
\node at (2.2,1.3) {Realistic};

% Amdahl's limit
\draw[thick, red, dashed] (0,0) .. controls (1,0.7) and (2,1.2) .. (2.8,1.3);
\node at (2.2,0.9) {Amdahl};

% Labels
\node at (1.5,3.5) {\textbf{Scaling Behavior}};
}
\end{center>
\end{column>
\end{columns>
\end{frame>

\begin{frame}{Applications and Use Cases}
\begin{columns}
\begin{column}{0.6\textwidth}
\concept{Suitable Applications}:
\begin{itemize}
\item \textbf{Scientific Computing}: Numerical simulations, modeling
\item \textbf{Image Processing}: Pixel-level operations, filtering
\item \textbf{Signal Processing}: Digital filters, transforms
\item \textbf{Machine Learning}: Matrix operations, neural networks
\end{itemize>

\vspace{0.3cm}
\concept{Application Characteristics}:
\begin{itemize}
\item Data-parallel computations
\item Regular memory access patterns
\item Minimal inter-processor communication
\item Embarrassingly parallel problems
\end{itemize>

\vspace{0.3cm}
\concept{Modern Implementations}:
\begin{itemize}
\item \textbf{GPUs}: Graphics processing units
\item \textbf{Vector Processors}: Specialized SIMD machines
\item \textbf{Systolic Arrays}: Regular computation patterns
\item \textbf{FPGA Arrays}: Reconfigurable parallel processing
\end{itemize>
\end{column>
\begin{column>{0.4\textwidth>
\begin{center>
\textbf{Application Domains}
\vspace{0.3cm}

\tikz{
% Application categories
\node at (1.5,3.5) {\textbf{N-OS Applications}};

\draw[thick] (0.5,3) rectangle (2.5,3.2);
\node at (1.5,3.1) {Scientific Computing};

\draw[thick] (0.5,2.6) rectangle (2.5,2.8);
\node at (1.5,2.7) {Image Processing};

\draw[thick] (0.5,2.2) rectangle (2.5,2.4);
\node at (1.5,2.3) {Signal Processing};

\draw[thick] (0.5,1.8) rectangle (2.5,2);
\node at (1.5,1.9) {Machine Learning};

% Characteristics
\node at (1.5,1.5) {\textbf{Common Traits}};
\node at (1.5,1.2) {Data Parallel};
\node at (1.5,0.9) {Regular Patterns};
\node at (1.5,0.6) {SIMD Friendly};
}
\end{center>
\end{column>
\end{columns>
\end{frame>

% Summary
\section{Summary}

\begin{frame}{Chapter Summary}
\begin{columns}
\begin{column}{0.6\textwidth}
\concept{Key Concepts Covered}:
\begin{itemize>
\item N-Order Systems extend sequential computing to parallel
\item Cellular Automata provide theoretical foundation
\item Controlled Cellular Automata enable practical implementation
\item SIMD paradigm: Single Instruction, Multiple Data
\item Space-time trade-offs optimize resource utilization
\item Scalability challenges and solutions
\end{itemize>

\vspace{0.3cm}
\concept{Architectural Innovations}:
\begin{itemize}
\item Distribution networks for command broadcast
\item Structured state spaces with register files
\item RALU execution units in each cell
\item Matrix-based state representation
\item Efficient data transfer mechanisms
\end{itemize>

\vspace{0.3cm}
\concept{Modern Relevance}:
\begin{itemize}
\item Foundation for GPU architectures
\item Basis for vector processing systems
\item Principles apply to many-core processors
\item Relevant to AI/ML accelerators
\end{itemize>
\end{column>
\begin{column>{0.4\textwidth>
\begin{center>
\textbf{N-OS Evolution}
\vspace{0.3cm}

\tikz{
\node[draw, circle, minimum width=1.5cm] (center) at (0,0) {N-OS};
\node[draw, rectangle] (ca) at (0,1.5) {Cellular Automata};
\node[draw, rectangle] (cca) at (-1.5,0.5) {Controlled CA};
\node[draw, rectangle] (simd) at (1.5,0.5) {SIMD};
\node[draw, rectangle] (modern) at (0,-1.5) {Modern GPUs};

\draw[->] (ca) -- (center);
\draw[->] (center) -- (cca);
\draw[->] (center) -- (simd);
\draw[->] (center) -- (modern);
}
\end{center>
\end{column>
\end{columns>
\end{frame>

\begin{frame}{Discussion Questions}
\begin{enumerate}
\item How do N-Order Systems address the limitations of sequential computing?
\item What are the key differences between theoretical and controlled cellular automata?
\item How does the distribution network design affect system scalability?
\item What types of problems are best suited for N-OS architectures?
\item How do space-time trade-offs influence system design decisions?
\item What are the main challenges in achieving linear speedup with parallel systems?
\item How do modern GPU architectures relate to the N-OS concepts?
\item What role does data locality play in N-OS performance?
\end{enumerate}
\end{frame>

\end{document}

